
\textbf{Benchmark design.} TruthfulQA (Lin et al., 2021) evaluates whether models generate truthful responses across 817 questions spanning 38 fine-grained categories, which technical reports typically aggregate into 6-12 broader categories. MMLU (Hendrycks et al., 2021) assesses knowledge across 14,042 questions in 57 subjects grouped into four domains (STEM, humanities, social sciences, other). Both benchmarks are widely used in model evaluation and appear in most major technical reports since their publication.

\textbf{Error analysis methods.} Existing approaches to understanding model failures typically require either full model access (for activation analysis or attention visualization) or the ability to run inference (for generating item-level responses to analyze). These methods produce detailed insights but create cost and access barriers. Our approach differs by examining only published aggregated results, eliminating the need for model access or re-evaluation.

\textbf{Meta-analysis of benchmarks.} Prior work comparing models typically aggregates scores at the model level (e.g., "Model X achieves Y\% on Benchmark Z") without systematically examining category-level structure within benchmarks. Some analyses extract performance trends across model generations, but these typically use overall scores rather than per-category breakdowns.

\textbf{Weak supervision.} The weak supervision paradigm (Ratner et al., 2019) demonstrates that machine learning can proceed with coarse-grained or noisy labels rather than precise ground truth. In this framework, category-level labels could serve as weak supervision for item-level analysis. However, this requires that category-level labels exist and are available. Our work validates this prerequisite for published benchmark data but does not test the effectiveness of weak supervision approaches that might build on it.
